\subsubsection{Wall jump a wall slide}

Další implementovanou mechanikou bylo skákání a klouzání po stěnách, tedy \textit{Wall jump} a \textit{Wall slide}. Tyto dvě funkce spolu úzce souvisejí. \textit{Wall jump} umožňuje hráči opětovně vyskočit při kontaktu se stěnou ze strany, čímž se otevírá možnost postupného pohybu mezi dvěma vertikálními plochami a dosažení vyšších částí úrovně.

\textit{Wall slide} doplňuje předchozí mechaniku tím, že upravuje vertikální rychlost při kontaktu se stěnou. Bez této úpravy by hráč podléhal standardní gravitaci a rychlosti a neměl by dostatečný časový prostor pro provedení dalšího skoku. \textit{Wall slide} proto zpomaluje pád podél stěny a umožňuje kontrolovanější vertikální pohyb. Díky tomu se dá teoreticky nastavit, aby hráč nemohl vůbec spadnout a byl \textit{přilepený} na stěně, což se k některým hrám může hodit, ovšem my jsme se rozhodli pro to, aby hráč stále padal i na stěně.